\documentclass[12pt]{article}
\usepackage[utf8]{inputenc}
\usepackage[T1]{fontenc}
\usepackage{hyperref}
\title{Banyjima Language Resources for K-12 Curriculum Development}
\author{Compiled by Claude Sonnet 5.5 (claude-sonnet-5-5) for Andr\'as Kornai}
\date{}
\begin{document}
\maketitle

\section{Introduction}
Banyjima (Panyjima; ISO 639-3 pnw) is a Pama-Nyungan language of the Pilbara region of Western Australia. Search results report 140 speakers in the 2021 census, while the Wangka Maya site gives an estimate of fewer than 50 speakers in 2004. The main resources come from the Wangka Maya Pilbara Aboriginal Language Centre.

\section{Primary Educational Resources}
\begin{itemize}
\item Wangka Maya Banyjima topical wordlist with accompanying audio and illustrations (2010), reported in search results; no direct link was seen.
\end{itemize}
No school curriculum or programme was found.

\section{Linguistic and Reference Materials}
\begin{itemize}
\item \emph{Banyjima Dictionary with English--Banyjima wordlist and topic wordlists}, compiled by Annie Edwards-Cameron, Wangka Maya (2018, 142 pages; the bookseller listing shows it sold out): \url{https://elizabethsbookshop.com.au/shop/australiana/west-australiana/banyjima-banyjima-dictionary-with-english-banyjima-wordlist-and-topic-wordlists}
\item \emph{Bilybara} (1992), a Pilbara reference book with a Banyjima spelling system, words and sentences, a word list and texts; described as not comprehensive: \url{https://nthieberger.net/Bilybara.pdf}
\item Wikipedia, Panyjima language: \url{https://en.wikipedia.org/wiki/Panyjima_language}
\item AIATSIS Austlang entry A53: \url{https://collection.aiatsis.gov.au/austlang/language/A53}
\end{itemize}
No dedicated grammar was identified in the results.

\section{Religious and Community Resources}
\begin{itemize}
\item Wangka Maya Pilbara Aboriginal Language Centre, Banyjima page: \url{https://www.wangkamaya.org.au/pilbara-languages/01-banyjima}
\end{itemize}

\section{Digital and Audio Resources}
Audio accompanies the 2010 topical wordlist according to the search summary, but no online audio link was seen. The Wangka Maya and AIATSIS pages above should be checked for recordings.

\section{Curriculum Development Context}
The dictionary and wordlists give a base for vocabulary-led teaching. A curriculum should be developed with Banyjima elders and Wangka Maya, who hold the materials and authority over them. The dictionary's availability needs checking with the publisher.
\end{document}